% Section 3.5: Cross-Correlations as a Complementary Probe
% Structure: Introduction -> The Angular Power Spectrum -> Statistical Sensitivity

\section{Cross-Correlations as a Complementary Probe}
\label{sec:cross_correlations}

% TODO: review in light of the xcorr chapter. are there some overlaps? maybe we can bring or bridge the xcorr formalism appendix here, and remove that appendix. 

The methods discussed so far in this chapter---frequentist and Bayesian inference, simulation-based approaches, machine learning, and domain adaptation---all operate on data from a single instrument, extracting information about source populations from the gamma-ray sky alone.
We now consider a qualitatively different approach: correlating the gamma-ray emission with an independent gravitational tracer of the dark matter distribution, such as a galaxy catalog.
The angular cross-power spectrum between these two observables provides sensitivity to the collective dark matter signal from unresolved cosmological structure, offering a complementary probe that goes beyond what any single dataset can achieve.

Cross-correlations are powerful precisely because they exploit the fact that dark matter, regardless of its particle nature, is the dominant source of gravitational clustering.
If dark matter annihilates or decays into gamma rays, the resulting emission must trace the same large-scale structure that is independently mapped by galaxy surveys.
By measuring the spatial correlation between the two fields, one can isolate the collective dark matter contribution from the aggregate of unresolved cosmological structure---a statistical signal that no individual source produces, and one that would be undetectable in either dataset alone.

We introduce the cross-correlation formalism in two steps.
Section~\ref{sec:angular_power_spectrum} defines the angular power spectrum, the Limber approximation that projects three-dimensional clustering into two-dimensional observables, and the window functions that encode the redshift distribution of each signal.
Section~\ref{sec:statistical_sensitivity} presents the statistical tools---the signal-to-noise ratio and the $\Delta\chi^2$ test statistic---used to forecast the sensitivity of a given experiment to dark matter.
The full application of this formalism, including the halo model computation and the detailed sensitivity forecasts for the Cherenkov Telescope Array Observatory, is deferred to Chapter~10.

\subsection{The Angular Power Spectrum}
\label{sec:angular_power_spectrum}

The observable sky maps produced by gamma-ray telescopes and galaxy surveys are two-dimensional projections of the three-dimensional distribution of matter and radiation in the Universe.
To study the spatial correlations between these projected fields, we expand them in the natural basis of the sphere: the spherical harmonics $Y_{\ell m}(\hat{n})$.
The angular power spectrum (APS) $C_\ell$ then quantifies the amplitude of these correlations as a function of the angular scale, parameterized by the multipole $\ell$.

We begin by considering a generic intensity field $I_g(\hat{n})$ along the direction $\hat{n}$, which can be written as an integral over the radial comoving distance $\chi(z)$ \cite{Fornengo:2013rga, Camera:2012cj}:
\begin{equation}
	I_g(\hat{n}) = \int d\chi\, g(\chi, \hat{n})\, \tilde{W}(\chi)\,,
	\label{eq:intensity_field}
\end{equation}
where $g(\chi, \hat{n})$ is the density field of the source and $\tilde{W}(\chi)$ is the window function.
By defining a normalized window function $W(\chi) = \langle g \rangle \tilde{W}(\chi)$, the mean intensity simplifies to $\langle I_g \rangle = \int d\chi\, W(\chi)$.
The intensity fluctuations $\delta I_g(\hat{n}) \equiv I_g(\hat{n}) - \langle I_g \rangle$ can then be expanded in spherical harmonics as
\begin{equation}
	\delta I_g(\hat{n}) = \langle I_g \rangle \sum_{\ell m} a_{\ell m}\, Y_{\ell m}(\hat{n})\,,
	\label{eq:spherical_harmonic_expansion}
\end{equation}
where the dimensionless coefficients $a_{\ell m}$ encode the angular structure of the fluctuation field \cite{Fornengo:2013rga}.

For two intensity fields labelled $i$ and $j$ (for instance, the gamma-ray flux and the galaxy number counts), the angular power spectrum is defined as the ensemble average of the product of their expansion coefficients \cite{Fornengo:2013rga, Camera:2012cj}:
\begin{equation}
	C_\ell^{ij} = \frac{1}{2\ell+1} \left\langle \sum_{m=-\ell}^{\ell} a_{\ell m}^{(i)}\, a_{\ell m}^{(j)*} \right\rangle\,.
	\label{eq:aps_definition}
\end{equation}
When $i = j$, the expression gives the auto-correlation APS, which measures the intrinsic clustering of a single field.
When $i \neq j$, we obtain the cross-correlation APS, which measures the spatial overlap between two distinct observables.
By construction, $C_\ell^{ij} = C_\ell^{ji}$.

To connect the measured two-dimensional APS to the underlying three-dimensional power spectrum $P_{ij}(k,\chi)$, which contains the physics of the source distribution, we apply the Limber approximation \cite{Limber:1953apj}.
This approximation, valid for multipoles $\ell \gg 1$, replaces the oscillatory integrals over spherical Bessel functions with a single line-of-sight integral by identifying the wavenumber $k = \ell / \chi$.
The result is a compact expression for the angular power spectrum \cite{Fornengo:2013rga, Camera:2012cj}:
\begin{equation}
	C_\ell^{ij} = \frac{1}{\langle I_i \rangle \langle I_j \rangle} \int \frac{d\chi}{\chi^2}\, W_i(\chi)\, W_j(\chi)\, P_{ij}\!
	\left(k = \frac{\ell}{\chi},\, \chi\right)\,.
	\label{eq:limber_aps}
\end{equation}
The amplitude of the cross-correlation signal depends on the overlap between the window functions $W_i(\chi)$ and $W_j(\chi)$ and on the three-dimensional cross-power spectrum at the relevant scales.
A significant cross-correlation therefore requires that both observables trace the same large-scale structure at similar redshifts \cite{Fornengo:2013rga}.
This idea of correlating a gamma-ray signal with a gravitational tracer of the dark matter distribution was pioneered by Camera et al.
\ \cite{Camera:2012cj}, who first showed that the cross-correlation of gamma-ray emission with weak-lensing cosmic shear provides a viable channel for indirect dark matter searches.

The window functions are the central ingredients of this formalism, as they encode how each observable is distributed along the line of sight.
Their specific form depends on the physical origin of the signal.
For a galaxy catalog acting as a gravitational tracer of the large-scale structure, the window function is simply the normalized redshift distribution of the surveyed galaxies, $W_g(z) = dN_g/dz$ \cite{Pinetti:2025hgd, Fornengo:2013rga}.
The angular power spectrum from annihilating dark matter was first formulated by Ando and Komatsu \cite{Ando:2005xg}, and the formalism was later extended to multi-wavelength and multi-tracer cross-correlations by Fornengo and Regis \cite{Fornengo:2013rga}.

For dark matter annihilation, the gamma-ray window function takes the form \cite{Fornengo:2013rga, Ando:2005xg}
\begin{equation}
	W_{\gamma,\mathrm{ann}}(E, z) = \frac{(\Omega_\mathrm{DM} \rho_c)^2}{4\pi}\, \frac{\langle \sigma v \rangle}{2 m_\chi^2}\, (1+z)^3\, \Delta^2(z)\, \frac{dN_a}{dE}[E(1+z)]\, e^{-\tau(E,z)}\,,
	\label{eq:window_dm_ann}
\end{equation}
where $\langle \sigma v \rangle$ is the velocity-averaged annihilation cross section, $m_\chi$ is the dark matter particle mass, $\Delta^2(z) = \langle \rho^2 \rangle / \bar{\rho}^2$ is the clumping factor accounting for the enhancement due to substructure within halos, $dN_a/dE$ is the photon spectrum per annihilation event, and $\tau(E,z)$ is the optical depth for absorption on the extragalactic background light.
For decaying dark matter, the window function scales linearly with the density rather than quadratically, and the clumping factor does not appear \cite{Fornengo:2013rga}.
The annihilation window function is strongly peaked at low redshifts, where the local dark matter density is highest and the clumping factor most effective, while astrophysical source populations (such as blazars) peak at higher redshifts.
This difference in redshift distribution is what makes the cross-correlation with low-redshift galaxy catalogs particularly powerful for isolating a dark matter signal \cite{Pinetti:2025hgd, Fornengo:2013rga}.

The three-dimensional power spectrum $P_{ij}(k,\chi)$ in the Limber integral describes the spatial clustering of the correlated fields.
It is computed within the halo model framework, which partitions all matter in the Universe into discrete, collapsed dark matter halos \cite{Cooray:2002dia}.
Under this assumption, the power spectrum decomposes into two contributions \cite{Cooray:2002dia, Fornengo:2013rga}:
\begin{equation}
	P_{ij}(k, z) = P_{ij}^{\mathrm{1h}}(k, z) + P_{ij}^{\mathrm{2h}}(k, z)\,.
	\label{eq:halo_model_decomposition}
\end{equation}
The one-halo term $P^{\mathrm{1h}}$ accounts for correlations between two points within the same individual halo; it probes the internal mass distribution of halos and dominates on small physical scales (large $k$).
The two-halo term $P^{\mathrm{2h}}$ captures correlations between points in two distinct halos; it traces the underlying linear matter power spectrum modulated by the halo bias factors and dominates on large scales (small $k$).
Both terms are expressed as integrals over the halo mass function $dn/dM$ and the Fourier-transformed density profiles of the relevant tracers \cite{Cooray:2002dia, Fornengo:2013rga}.
The detailed computation of these terms for the specific gamma-ray and galaxy populations considered in this thesis is deferred to Chapter~10, where the full cross-correlation analysis is presented.

\subsection{Statistical Sensitivity}
\label{sec:statistical_sensitivity}

Given the cross-correlation angular power spectrum $C_\ell^{\gamma g}$ defined in the previous section, we now turn to the statistical tools used to assess how well a given experiment can detect or constrain a dark matter signal.
Two quantities are central to this assessment: the signal-to-noise ratio (SNR) and the $\Delta\chi^2$ test statistic.

The detectability of the cross-correlation signal from a given source population (whether astrophysical or dark matter) is quantified by the signal-to-noise ratio \cite{Pinetti:2025hgd}:
\begin{equation}
	\mathrm{SNR} = \sqrt{\sum_{\ell,\, i} \left(\frac{C_{\ell,i}^{\gamma g}}{\Delta C_{\ell,i}^{\gamma g}}\right)^2}\,,
	\label{eq:snr}
\end{equation}
where the sum runs over the multipoles $\ell$ included in the analysis and the energy bins $i$. The variance of the cross-correlation spectrum, assuming Gaussian-distributed fluctuations, is given by \cite{Pinetti:2025hgd}
\begin{equation}
	(\Delta C_\ell^{\gamma g})^2 = \frac{1}{(2\ell+1)\, f_\mathrm{sky}} \left[(C_\ell^{\gamma g})^2 + \left(C_\ell^{\gamma\gamma} + \frac{C_N}{B_\ell^2}\right)(C_\ell^{gg} + C_{N_g})\right]\,,
	\label{eq:variance_cross}
\end{equation}
where $f_\mathrm{sky}$ is the observed sky fraction, $C_\ell^{\gamma\gamma}$ and $C_\ell^{gg}$ are the gamma-ray and galaxy auto-correlation spectra, $C_N$ is the photon noise of the gamma-ray detector, $C_{N_g}$ is the shot noise of the galaxy catalog, and $B_\ell$ is the beam window function encoding the finite angular resolution (point-spread function) of the instrument.
No beam correction enters the galaxy term, as the angular resolution of galaxy catalogs (typically arcseconds) is negligible at the scales relevant for this analysis \cite{Pinetti:2025hgd}.

The structure of the variance formula reveals a fundamental advantage of cross-correlation over auto-correlation for dark matter searches.
The corresponding auto-correlation variance is \cite{Pinetti:2025hgd}
\begin{equation}
	(\Delta C_\ell^{\gamma\gamma})^2 = \frac{2}{(2\ell+1)\, f_\mathrm{sky}} \left(C_\ell^{\gamma\gamma} + \frac{C_N}{B_\ell^2}\right)^2\,,
	\label{eq:variance_auto}
\end{equation}
where the photon noise $C_N$ now enters quadratically and adds directly to the signal.
In an auto-correlation measurement, $C_\ell^{\gamma\gamma} + C_N$ must be measured and the noise accurately estimated and subtracted to recover the true clustering spectrum.
In cross-correlation, by contrast, the noise sources in the gamma-ray map and in the galaxy catalog are physically independent: their cross-noise vanishes on average.
The photon noise $C_N$ and galaxy shot noise $C_{N_g}$ therefore contribute only to the statistical uncertainty (the error bars) of the cross-correlation measurement, not to the signal itself.
This allows faint correlated signals---such as a subdominant dark matter contribution---to be extracted even in noisy datasets where the auto-correlation would be completely dominated by instrumental noise \cite{Pinetti:2025hgd, Camera:2012cj}.

To assess the sensitivity to a dark matter signal specifically, we compare two hypotheses by means of a $\Delta\chi^2$ test statistic \cite{Pinetti:2025hgd}:
\begin{equation}
	\Delta\chi^2 = \sum_{\ell,\, i} \left(\frac{C_{\ell,i}^{\gamma_\mathrm{astro+DM}\, g}}{\Delta C_{\ell,i}^{\gamma_\mathrm{astro+DM}\, g}}\right)^2 - \sum_{\ell,\, i} \left(\frac{C_{\ell,i}^{\gamma_\mathrm{astro}\, g}}{\Delta C_{\ell,i}^{\gamma_\mathrm{astro}\, g}}\right)^2\,,
	\label{eq:delta_chi2}
\end{equation}
where $C_{\ell}^{\gamma_\mathrm{astro}\, g}$ is the predicted cross-correlation signal under the null hypothesis (astrophysical sources only) and $C_{\ell}^{\gamma_\mathrm{astro+DM}\, g}$ includes an additional dark matter component.
The test statistic measures how much the inclusion of a dark matter signal improves the fit relative to the astrophysical-only model.
In Chapter~\ref{ch:8}, this formalism is applied to forecast the sensitivity of the Cherenkov Telescope Array Observatory (CTAO) to dark matter annihilation and decay, by cross-correlating the predicted gamma-ray emission with the 2MASS galaxy catalog \cite{Pinetti:2025hgd}.
